% Figure from MATH 493 notes, section 31. Compile with the notes preamble.
\begin{tikzpicture}[el/.style={font=\scriptsize, inner sep=0pt}, y=0.4cm]
  % columns: x-centre / half-width / cycle type / size / elements (top to bottom)
  \foreach \x/\w/\t/\s/\els in {
      0/0.4/1^4/1/{e},
      1.25/0.7/2^2/3/{(1\,2)(3\,4), (1\,3)(2\,4), (1\,4)(2\,3)},
      2.75/0.6/3\,1/8/{(1\,2\,3), (1\,3\,2), (1\,2\,4), (1\,4\,2), (1\,3\,4), (1\,4\,3), (2\,3\,4), (2\,4\,3)},
      4.1/0.6/2\,1^2/6/{(1\,2), (1\,3), (1\,4), (2\,3), (2\,4), (3\,4)},
      5.45/0.6/4/6/{(1\,2\,3\,4), (1\,2\,4\,3), (1\,3\,2\,4), (1\,3\,4\,2), (1\,4\,2\,3), (1\,4\,3\,2)}} {
    \node[font=\scriptsize, black!70] at (\x,0.8) {$\t$};
    \node[font=\scriptsize, figR] at (\x,-9.4) {$\s$};
    \foreach \e [count=\k] in \els { \node[el] at (\x,-\k+0.3) {$\e$}; \xdef\last{\k} }
    \begin{scope}[on background layer]
      \draw[figB, thick, rounded corners=4pt, fill=figB!8] (\x-\w,0.1) rectangle (\x+\w,-\last-0.3);
    \end{scope}
  }
  \node[font=\scriptsize, black!70, anchor=east] at (-0.55,0.8) {cycle type};
  \node[font=\scriptsize, figR, anchor=east] at (-0.55,-9.4) {class size};
\end{tikzpicture}
